%---------------------------------------------------------------------------%
%->> Backmatter
%---------------------------------------------------------------------------%
\chapter[致谢]{致\quad 谢}\chaptermark{致\quad 谢}% syntax: \chapter[目录]{标题}\chaptermark{页眉}
%\thispagestyle{noheaderstyle}% 如果需要移除当前页的页眉
%\pagestyle{noheaderstyle}% 如果需要移除整章的页眉

此处填写致谢。


\rightline{2023年6月}
\chapter{作者简历及攻读学位期间发表的学术论文与其他相关学术成果}

\section*{作者简历：}
2016年9月——2020年6月，在复旦大学计算机科学技术学院获得学士学位。

2020年9月——2022年6月，在中国科学院计算技术研究所获得硕士学位。

2022年9月——至今，在中国科学院计算技术研究所攻读博士学位（硕博连读）。

工作经历：无。% 待确认：如无工作经历可保留“无”，否则按“起止年月，单位，职务”填写

\section*{已发表（或正式接受）的学术论文：}

{
\setlist[enumerate]{}% restore default behavior
\begin{enumerate}[nosep]
    \item 李润宽, 胡斌, 刘晓东, 赵晓芳. I-KASLR：面向智能电网等实时关键系统的内置内核代码持续重随机化方法[J]. 高技术通讯, 2026, 36(1): 1-14.
    \item 李润宽, 胡斌, 赵晓芳, 蒋德钧. 基于编译的 Linux 内核地址空间布局控制方法[J]. 高技术通讯, 2025, 34(9): 921-934. DOI: 10.3772/j.issn.1002-0470.2025.09.002.
    \item Zihao Chang, \textbf{Runkuan Li}, JiHan Lin. Chaos: Function Granularity Runtime Address Layout Space Randomization for Kernel Module[C]. // Proceedings of the 15th ACM SIGOPS Asia-Pacific Workshop on Systems (APSys 2024), 2024.% 待补：页码与 DOI
\end{enumerate}
}

\section*{申请或已获得的专利：}

{
\setlist[enumerate]{}% restore default behavior
\begin{enumerate}[nosep]
    \item 李润宽, 胡斌, 王盈, 刘洁, 常子豪, 张焕, 林纪涵, 蒋德钧, 王卅, 赵晓芳. 一种基于编译器的细粒度程序执行流随机化控制方法: 中国, 202410206829.X[P]. 实质审查中.% 待确认：申请日与当前法律状态
    \item 李润宽, 刘洁, 胡斌, 蒋德钧, 赵晓芳. 一种基于隔离机制的 Linux 内核代码的地址空间布局触发式随机化方法: 中国, 202510342155.3[P]. 实质审查中.% 待确认：申请日与当前法律状态
\end{enumerate}
}

\section*{参加的研究项目及获奖情况：}

% 待填：参加的研究项目（项目名称、编号、来源、起止年月、本人承担的工作）与获奖情况；若均无，此节可删去


\cleardoublepage[plain]% 让文档总是结束于偶数页，可根据需要设定页眉页脚样式，如 [noheaderstyle]
%---------------------------------------------------------------------------%
